% =================================================
% Домашняя работа: решения для учителя — скорректированная версия
% =================================================
\newpage
\homeworktitle
  {Решения домашней работы}
  {Ответы, ключевые шаги и диагностические комментарии}

\begin{teacheranswer}[Задание 1]
  а)
  \[
    \lg^3 A=(\lg A)^3,
    \qquad
    \lg(A^3)\text{ — логарифм аргумента }A^3,
  \]
  \[
    \log_3 A\text{ — логарифм }A\text{ по основанию }3.
  \]
  б)
  \[
    \lg^2 1000=3^2=9,
    \qquad
    \lg(1000^2)=\lg10^6=6.
  \]
  в) В первой записи возводится в квадрат значение функции, во второй —
  изменяется аргумент функции.
\end{teacheranswer}

\begin{criterionbox}
  Проверяется чтение функциональной нотации, а не только вычислительный навык.
\end{criterionbox}

\begin{teacheranswer}[Задание 2]
  Левая часть определена при \(x\ne1\), правая — только при \(x>1\).
  Равенство, верное на всей ОДЗ левой части:
  \[
    \lg((x-1)^2)=2\lg|x-1|,
    \qquad x\ne1.
  \]
\end{teacheranswer}

\begin{teacheranswer}[Задание 3]
  Пусть \(t=y^2\), \(t\ge0\). Тогда
  \[
    t^2-t-12\ge0,
    \qquad
    (t-4)(t+3)\ge0.
  \]
  Отсюда \(t\le-3\) или \(t\ge4\). С учётом \(t\ge0\): \(t\ge4\).
  Поэтому
  \[
    y^2\ge4
    \quad\Longleftrightarrow\quad
    y\le-2\ \text{или}\ y\ge2.
  \]
  \textbf{Ответ:} \(( -\infty,-2]\cup[2,+\infty)\).
\end{teacheranswer}

\begin{teacheranswer}[Задание 4]
  ОДЗ: \(x\ne1\). Положим
  \[
    y=\log_2|x-1|.
  \]
  Тогда
  \[
    \log_2((x-1)^2)=2y,
    \qquad
    \log_2((x-1)^4)=4y.
  \]
  Получаем
  \[
    (2y)^4-(4y)^2\ge192
    \quad\Longleftrightarrow\quad
    y^4-y^2-12\ge0.
  \]
  По заданию 3:
  \[
    y\le-2\quad\text{или}\quad y\ge2.
  \]
  Следовательно,
  \[
    0<|x-1|\le2^{-2}=\frac14
    \quad\text{или}\quad
    |x-1|\ge2^2=4.
  \]
  Первая ветвь даёт
  \[
    x\in\left[\frac34,1\right)\cup\left(1,\frac54\right].
  \]
  Вторая ветвь даёт
  \[
    x\in(-\infty,-3]\cup[5,+\infty).
  \]
  \textbf{Ответ:}
  \[
    (-\infty,-3]
    \cup\left[\frac34,1\right)
    \cup\left(1,\frac54\right]
    \cup[5,+\infty).
  \]
\end{teacheranswer}

\begin{criterionbox}
  Проверяются: модуль, две замены, две ветви для логарифма и исключение \(x=1\).
\end{criterionbox}

\begin{teacheranswer}[Задание 5]
  Подходит вариант A:
  \[
    t=\lg^2|x^2-1|
    =\bigl(\lg|x^2-1|\bigr)^2.
  \]
  После вынесения коэффициента \(16\) исходное выражение становится
  \[
    16(t^2-t).
  \]
  Вариант Б использует основание \(2\), а вариант В не является квадратом
  значения логарифма.
\end{teacheranswer}

\begin{teacheranswer}[Задание 6]
  ОДЗ: \(x\ne-2\). Положим
  \[
    y=\log_3|x+2|.
  \]
  Получаем \(y^4-y^2-12\ge0\), откуда
  \[
    y\le-2\quad\text{или}\quad y\ge2.
  \]
  Следовательно,
  \[
    0<|x+2|\le3^{-2}=\frac19
    \quad\text{или}\quad
    |x+2|\ge3^2=9.
  \]
  Первая ветвь:
  \[
    x\in\left[-\frac{19}{9},-2\right)
    \cup\left(-2,-\frac{17}{9}\right].
  \]
  Вторая ветвь:
  \[
    x\in(-\infty,-11]\cup[7,+\infty).
  \]
  \textbf{Ответ:}
  \[
    (-\infty,-11]
    \cup\left[-\frac{19}{9},-2\right)
    \cup\left(-2,-\frac{17}{9}\right]
    \cup[7,+\infty).
  \]
\end{teacheranswer}

\begin{teacheranswer}[Задание 7]
  ОДЗ: \(x\ne\pm1\). Положим
  \[
    y=\log_2|x^2-1|.
  \]
  После двух замен получаем
  \[
    y\le-2\quad\text{или}\quad y\ge2,
  \]
  поэтому
  \[
    0<|x^2-1|\le\frac14
    \quad\text{или}\quad
    |x^2-1|\ge4.
  \]
  Первая ветвь даёт
  \[
    \frac34\le x^2<1
    \quad\text{или}\quad
    1<x^2\le\frac54.
  \]
  Вторая ветвь сводится к \(x^2\ge5\), поскольку \(x^2\le-3\) невозможно.
  \textbf{Ответ:}
  \[
    \begin{aligned}
      x\in{}&(-\infty,-\sqrt5]
      \cup\left[-\frac{\sqrt5}{2},-1\right)
      \cup\left(-1,-\frac{\sqrt3}{2}\right]\\
      &\cup\left[\frac{\sqrt3}{2},1\right)
      \cup\left(1,\frac{\sqrt5}{2}\right]
      \cup[\sqrt5,+\infty).
    \end{aligned}
  \]
\end{teacheranswer}

\begin{teacheranswer}[Задание 8]
  Поскольку \(A\ne0\), положим
  \[
    y=\log_b|A|.
  \]
  Тогда
  \[
    \log_b(A^2)=2y,
    \qquad
    \log_b(A^4)=4y.
  \]
  Получаем \(y^4-y^2-12\ge0\), поэтому
  \[
    y\le-2\quad\text{или}\quad y\ge2.
  \]
  При \(b>1\):
  \[
    0<|A|\le b^{-2}
    \quad\text{или}\quad
    |A|\ge b^2.
  \]
  \textbf{Ответ:}
  \[
    A\in(-\infty,-b^2]
    \cup[-b^{-2},0)
    \cup(0,b^{-2}]
    \cup[b^2,+\infty).
  \]
\end{teacheranswer}

\begin{criterionbox}[Итоговая диагностика]
  Задания 1--2 проверяют чтение нотации и равносильность формул; задание 3 —
  алгебраическое ядро; задание 4 — сборку метода с опорой; задание 5 — выбор
  замены; задание 6 — перенос с минимальной опорой; задание 7 — полностью
  самостоятельное решение; задание 8 — обобщение.
\end{criterionbox}
